\documentclass{beamer}
\usepackage{graphicx}
\usepackage{xcolor}
\usepackage{tikz}
\usepackage{amsmath}

% 自定义颜色
\definecolor{purpleblue}{RGB}{80, 80, 180}
\definecolor{lightgray}{RGB}{230, 230, 230}
\definecolor{novelPurple}{RGB}{98, 54, 255}
\definecolor{darkgray}{RGB}{50, 50, 50}
\definecolor{softPurple}{RGB}{180, 170, 255}
\definecolor{softGray}{RGB}{235, 235, 245}
\definecolor{lightText}{RGB}{90, 90, 110}
\graphicspath{{static/course_materials/}{./}}

% 主题设置
\usetheme{Madrid}
\usecolortheme{dolphin}

\setbeamerfont{title}{size=\Large,series=\bfseries}
\setbeamerfont{frametitle}{size=\LARGE,series=\bfseries}
\setbeamercolor{frametitle}{fg=white, bg=novelPurple}
\setbeamercolor{title}{fg=white, bg=purpleblue}

\title[SAT Unit 3.2]{\textbf{SAT Unit 3.2 \\Percentages}}
\author{\textbf{Jack Zeng}}
\institute{\textcolor{white}{\textbf{Novel Prep}}}
\date{\today}

\begin{document}


% --- Title Slide ---
\begin{frame}
    \titlepage
    \vfill
    \begin{flushright}
        \textcolor{purpleblue}{\Huge \textbf{Novel Prep}}
    \end{flushright}
\end{frame}

\begin{frame}
    \begin{tikzpicture}[remember picture, overlay]
        \shade[top color=softPurple, bottom color=softGray]
            (current page.north west) rectangle (current page.south east);
        \fill[white, opacity=0.3] (current page.center) circle (4cm);
        \node[align=center] at (current page.center) {
            {\Huge\bfseries\textcolor{purpleblue}{Novel Prep}} \\[0.5cm]
            {\Large\itshape\textcolor{lightText}{Excellence in SAT Prep}}
        };
        \foreach \x in {1,2,3,4,5} {
            \fill[purpleblue, opacity=0.2] (rand*10-5, rand*7-3.5) circle (0.1+\x*0.05);
        }
    \end{tikzpicture}
\end{frame}

\begin{frame}{Overview}
    \tableofcontents
\end{frame}

\begin{frame}
\section{Introduction to Percent}
    \begin{tikzpicture}[remember picture, overlay]

        \shade[top color=softPurple, bottom color=softGray] 
            (current page.north west) rectangle (current page.south east);

        \fill[white, opacity=0.3] (current page.center) circle (4cm);
        

        \node[align=center] at (current page.center) {
            {\LARGE\bfseries\textcolor{purpleblue}{Introduction to Percent}} \\[0.5cm]
            {\Large\itshape\textcolor{lightText}{Excellence in SAT Prep}}
        };
        \foreach \x in {1,2,3,4,5} {
            \fill[purpleblue, opacity=0.2] (rand*10-5, rand*7-3.5) circle (0.1+\x*0.05);
        }
        ;
    \end{tikzpicture}
\end{frame}

% Slide 2: Introduction to Percent
\begin{frame}{Introduction to Percent}
    \textbf{Definition:} Percent means "per hundred." It represents a ratio of a number to 100.
    \[
    \text{Percent} = \frac{\text{Part}}{\text{Whole}} \times 100
    \]
    \textbf{Example:} 25 out of 100 is written as 25\%.
\end{frame}

% Slide 3: Percent Change Formula
\begin{frame}{Percent Change Formula}
    \textbf{Formula:}
    \[
    \text{Percent Change} = \frac{\text{New Value} - \text{Original Value}}{\text{Original Value}} \times 100
    \]
    \textbf{Example:} If a price increases from \$80 to \$90:
    \[
    \frac{90 - 80}{80} \times 100 = 12.5\%
    \]
    \textbf{Key Concept:} Positive percent change indicates an increase, negative indicates a decrease.
\end{frame}


\begin{frame}
\section{Percent Greater or Less Than}
    \begin{tikzpicture}[remember picture, overlay]

        \shade[top color=softPurple, bottom color=softGray] 
            (current page.north west) rectangle (current page.south east);

        \fill[white, opacity=0.3] (current page.center) circle (4cm);
        

        \node[align=center] at (current page.center) {
            {\LARGE\bfseries\textcolor{purpleblue}{Percent Greater or Less Than}} \\[0.5cm]
            {\Large\itshape\textcolor{lightText}{Excellence in SAT Prep}}
        };
        \foreach \x in {1,2,3,4,5} {
            \fill[purpleblue, opacity=0.2] (rand*10-5, rand*7-3.5) circle (0.1+\x*0.05);
        }
        ;
    \end{tikzpicture}
\end{frame}


% Slide 4: Percent Greater or Less Than
\begin{frame}{Percent Greater or Less Than}

\begin{center}
\begin{tabular}{|l|c|c|}
        \hline
        & \textbf{Translates to...} & \textbf{Example} \\
        \hline
        \(p\%\) of \(x\) & \(\dfrac{p}{100} \times x\) & \(15\%\) of \(x \Rightarrow \dfrac{15}{100} \times x = 0.15x\) \\
        \hline
        \(p\%\) greater than \(x\) & \(\left(1 + \dfrac{p}{100} \right)x\) & \(15\%\) greater than \(x \Rightarrow \left(1 + \dfrac{15}{100} \right)x = 1.15x\) \\
        \hline
        \(p\%\) less than \(x\) & \(\left(1 - \dfrac{p}{100} \right)x\) & \(15\%\) less than \(x \Rightarrow \left(1 - \dfrac{15}{100} \right)x = 0.85x\) \\
        \hline
    \end{tabular}
\end{center}


    \textbf{Example:} If 200 is \textbf{x}\% greater than 50, what is \textbf{x}?
    \[
   50 \times \left(1 + \frac{x}{100}\right) = 200
    \]
    \textbf{Key Takeaway:} If something is \textbf{50\% greater than} a value, multiply by 1.5.
\end{frame}

\begin{frame}{\textbf{Question: Percent Relationship Between \(a\), \(b\), and \(c\)}}
    The number \( a \) is \textbf{70\% less} than the positive number \( b \).

    The number \( c \) is \textbf{60\% greater} than \( a \).

    The number \( c \) is how many times \( b \)?
\end{frame}

\begin{frame}{\textbf{Answer Explanation: Percent Relationship Between \(a\), \(b\), and \(c\)}}
    \begin{itemize}
        \item Since \( a \) is 70\% less than \( b \), we express it as:
        \[
        a = (1 - 0.7)b = 0.3b
        \]
        \item Since \( c \) is 60\% greater than \( a \), we express it as:
        \[
        c = (1 + 0.6)a = 1.6a
        \]
        \item Substituting \( a = 0.3b \):
        \[
        c = 1.6(0.3b) = 0.48b
        \]
        \item Thus, the number \( c \) is \(\mathbf{0.48}\) times \( b \).
    \end{itemize}

    \textbf{Correct Answer:} \( \boxed{0.48} \)
\end{frame}

% Slide 7: Summary and Key Takeaways
\begin{frame}{Key Takeaways}
    \begin{itemize}
        \item Percent represents "per hundred."
        \item Use the \textbf{percent change formula} to calculate increases and decreases.
        \item "Percent greater than" and "percent less than" refer to the original value.
        \item Use division to find the original value when given the increased price.
    \end{itemize}
\end{frame}


\begin{frame}{\textbf{Question: Store Cost from Sale Price}}
The regular price of a shirt at a store is \(\$11.70\).

The sale price of the shirt is \textbf{80\% less} than the regular price,
and the sale price is \textbf{30\% greater} than the store's cost for the shirt.

What was the store's cost, in dollars, for the shirt?

\textbf{(Disregard the \$ sign when entering your answer.)}
\end{frame}

% Slide 3: Answer
\begin{frame}{\textbf{Answer Explanation: Store Cost from Sale Price}}
    \begin{itemize}
        \item \textbf{Step 1: Find the sale price of the shirt}
        \[
        \text{Sale Price} = (1 - 0.80) \times 11.70 = 0.20 \times 11.70 = 2.34
        \]
        \item \textbf{Step 2: Set up equation for the store’s cost}
        \[
        \text{Sale Price} = (1 + 0.30) \times \text{Cost}
        \]
        \[
        2.34 = 1.30 \times \text{Cost}
        \]
        \item \textbf{Step 3: Solve for Cost}
        \[
        \text{Cost} = \frac{2.34}{1.30} = 1.80
        \]
    \end{itemize}

    \textbf{Correct Answer:} \( \boxed{1.80} \)
\end{frame}
\begin{frame}
\section{Increased By vs. Decreased By}
    \begin{tikzpicture}[remember picture, overlay]

        \shade[top color=softPurple, bottom color=softGray] 
            (current page.north west) rectangle (current page.south east);

        \fill[white, opacity=0.3] (current page.center) circle (4cm);
        

        \node[align=center] at (current page.center) {
            {\LARGE\bfseries\textcolor{purpleblue}{Increased By vs. Decreased By}} \\[0.5cm]
            {\Large\itshape\textcolor{lightText}{Excellence in SAT Prep}}
        };
        \foreach \x in {1,2,3,4,5} {
            \fill[purpleblue, opacity=0.2] (rand*10-5, rand*7-3.5) circle (0.1+\x*0.05);
        }
        ;
    \end{tikzpicture}
\end{frame}
\begin{frame}{ ``Increased By'' vs. ``Decreased By''}
    
\textbf{Definitions:}
\begin{itemize}
    \item \textbf{Increased by \(p\%\):} Multiply by \(\left(1 + \frac{p}{100}\right)\)
    \item \textbf{Decreased by \(p\%\):} Multiply by \(\left(1 - \frac{p}{100}\right)\)
\end{itemize}

\textbf{Examples:}
\begin{itemize}
    \item \(120\) increased by \(25\%\): 
    \[
    120 \times \left(1 + \frac{25}{100}\right) = 120 \times 1.25 = 150
    \]
    \item \(150\) decreased by \(20\%\): 
    \[
    150 \times \left(1 - \frac{20}{100}\right) = 150 \times 0.8 = 120
    \]
\end{itemize}

\end{frame}

\begin{frame}{Important Notes}

\vspace{0.1cm}
\textbf{Important Notes:}
\begin{itemize}
    \item \textbf{Be careful:} A 20\% decrease followed by a 20\% increase does \textbf{not} return to the original value!
    \[
    100 \to 80 \to 96
    \]
    \item Always treat \textbf{``increased by'' and ``decreased by''} as multiplicative — not additive.
    \item If you are given a final value and percentage, \textbf{reverse the multiplier} to solve for the original.
\end{itemize}

\vspace{0.2cm}
\textbf{Key Insight:} Think in terms of \textbf{scaling}, not just adding or subtracting!
    
\end{frame}

\begin{frame}{\textbf{Question: Net Percentage Increase}}
The value of a collectible comic book increased by \(167\%\) from the end of \textbf{2011} to the end of \textbf{2012}, and then decreased by \(16\%\) from the end of \textbf{2012} to the end of \textbf{2013}. 

What was the net percentage increase in the value of the collectible comic book from the end of \textbf{2011} to the end of \textbf{2013}?

\begin{itemize}
    \item[\textbf{A.}] \(124.28\%\)
    \item[\textbf{B.}] \(140.28\%\)
    \item[\textbf{C.}] \(151.00\%\)
    \item[\textbf{D.}] \(209.72\%\)
\end{itemize}
\end{frame}


\begin{frame}{\textbf{Answer Explanation: Net Percentage Increase}}
Let the original value at the end of 2011 be \(100\) (for simplicity).

\begin{itemize}
    \item After a \(167\%\) increase: 
    \[
    100 + 1.67 \times 100 = 267
    \]
    \item After a \(16\%\) decrease: 
    \[
    267 \times (1 - 0.16) = 267 \times 0.84 = 224.28
    \]
    \item Net increase from original \(100\): 
    \[
    224.28 - 100 = 124.28
    \]
    The net increase is \(124.28\%\).
\end{itemize}

\textbf{Correct Answer:} \( \boxed{A} \)
\end{frame}


\begin{frame}
\section{Practice}
    \begin{tikzpicture}[remember picture, overlay]

        \shade[top color=softPurple, bottom color=softGray] 
            (current page.north west) rectangle (current page.south east);

        \fill[white, opacity=0.3] (current page.center) circle (4cm);
        

        \node[align=center] at (current page.center) {
            {\LARGE\bfseries\textcolor{purpleblue}{Practice}} \\[0.5cm]
            {\Large\itshape\textcolor{lightText}{Excellence in SAT Prep}}
        };
        \foreach \x in {1,2,3,4,5} {
            \fill[purpleblue, opacity=0.2] (rand*10-5, rand*7-3.5) circle (0.1+\x*0.05);
        }
        ;
    \end{tikzpicture}
\end{frame}


\begin{frame}{\textbf{Question: Percent Greater Relationship}}
In 2010, the number of houses built in Town A was \(n\) percent greater than the number of houses built in Town B. If \(700\) houses were built in Town A in 2010, which expression represents the number of houses built in Town B in 2010?
\begin{itemize}
    \item[\textbf{A.}] \(700 \left( \dfrac{100 + n}{100} \right)\)
    \item[\textbf{B.}] \(700 \left( \dfrac{100}{100 + n} \right)\)
    \item[\textbf{C.}] \(700 \left( \dfrac{100 - n}{100} \right)\)
    \item[\textbf{D.}] \(700 \left( \dfrac{100}{100 - n} \right)\)
\end{itemize}
\end{frame}


\begin{frame}{\textbf{Answer Explanation: Percent Greater Relationship}}
Let \(x\) be the number of houses in Town B.

Town A had \(n\%\) more than Town B, so:
\[
700 = x \left( \dfrac{100 + n}{100} \right)
\quad \Rightarrow \quad
x = 700 \cdot \left( \dfrac{100}{100 + n} \right)
\]

\textbf{Correct Answer:} \( \boxed{B} \)
\end{frame}


\begin{frame}{\textbf{Question: Percent of \(c\) that is \(a\)}}
The sum of positive numbers \(a\) and \(b\) is \(4.76\) times the positive number \(c\), and \(a\) is \(124\%\) of \(b\). What percent of \(c\) is \(a\)?
\begin{itemize}
    \item[\textbf{A.}] \(216.85\%\)
    \item[\textbf{B.}] \(238.00\%\)
    \item[\textbf{C.}] \(263.50\%\)
    \item[\textbf{D.}] \(295.12\%\)
\end{itemize}
\end{frame}


\begin{frame}{\textbf{Answer Explanation: Percent of \(c\) that is \(a\)}}
We are given:
\[
a + b = 4.76c \quad \text{and} \quad a = 1.24b
\]

Substitute \(a\) into the sum:
\[
1.24b + b = 4.76c \Rightarrow 2.24b = 4.76c \Rightarrow b = \dfrac{4.76c}{2.24}
\]

Now substitute \(b\) into \(a = 1.24b\):
\[
a = 1.24 \cdot \dfrac{4.76c}{2.24} = \dfrac{1.24 \cdot 4.76}{2.24}c = 2.635c
\]

Convert to percent:
\[
a = 263.5\% \text{ of } c
\]

\textbf{Correct Answer:} \( \boxed{C} \)
\end{frame}


\begin{frame}{\textbf{Question: Percent of \(z\) that is \(x\)}}
The sum of positive numbers \(x\) and \(y\) is \(3.5\) times the positive number \(z\), and \(x\) is \(160\%\) of \(y\). What percent of \(z\) is \(x\)?
\begin{itemize}
    \item[\textbf{A.}] \(212.00\%\)
    \item[\textbf{B.}] \(233.33\%\)
    \item[\textbf{C.}] \(256.00\%\)
    \item[\textbf{D.}] \(280.00\%\)
\end{itemize}
\end{frame}


\begin{frame}{\textbf{Answer Explanation: Percent of \(z\) that is \(x\)}}
\end{frame}

\begin{frame}{}
    \begin{tikzpicture}[remember picture, overlay]
        \shade[top color=novelPurple!40, bottom color=white]
            (current page.north west) rectangle (current page.south east);
        \fill[white, opacity=0.15] (current page.center) circle (5cm);
        \foreach \x in {1,2,3,4,5} {
            \fill[novelPurple, opacity=0.12] (rand*10-5, rand*7-3.5) circle (0.1+\x*0.05);
        }
    \end{tikzpicture}
    \centering
    \vspace{5em}
    {\Huge \textbf{\textcolor{novelPurple}{Thank You!}}} \\[1em]
    {\large \textcolor{gray}{We appreciate your time and focus.}} \\[3em]
    {\LARGE \textbf{\textcolor{novelPurple}{\textsf{Novel Prep}}}}
    \vfill
\end{frame}

\end{document}
